\scanpage{109}
\mspage{96}
% archive sheet number: 108

(5.8.13) on a $u^{*}(p')\in F$, on a $u^{*}(p')\in TF$,
i.e.
\[
 p'\in(u^{*})^{-1}(TF)=TF'\unsure{6|,}(\overset{\mathrm{déf}}=C'_{*}).
\]
la \deleted{dernière} égalité \deleted{\unsure{étant}} $(u^{*})^{-1}(\underset{\deleted{C_{*}}}{\underline{TF}})=TF'$,
\deleted{résultant de l'adjonction}
i.e.
\[
 (u^{*})^{-1}(C_{*})=(u_{!}(C))_{*}\ \unsure{\times|,}
\]
résultant de l'adjonction de $u_{!},u^{*}$. Au
total, on a donc
\[
 p'\in TF'\unsure{\times|,} (=C'_{*})\tag*{(5.8.16)}
\]
et comme $f'\in C'$, le diagramme $(\alpha)$
p.~114 nous montre qu'il existe
\[
 q':Y'\mathrel{\unsure{-|\longrightarrow}}X'_{1}\unsure{.,|,}
\]
avec $p'q'=\mathrm{id}$, $q'f'=i'$. Donc $f'$ est
facteur direct de $i'\in u_{!}(TC)^{\sim}$, donc
on a $f'\in u_{!}(TC)^{\sim}$, d'où la relation
\[
 TC'\deleted{\unsure{\cap W'}}\subset u_{!}(TC)^{\sim},
\]
ce qui achève de prouver que $(W',C',F')$
est un triple de Quillen spécial et clos.

Les deux relations (5.8.4) ont été prouvées
chemin faisant (la deuxième étant essentiellement triviale, et la première représentant
la difficulté principale \unsure{dans|de} la démonstration\unsure{.|)}
Reste \unsure{à établir|'établir} (5.8.5.). Mais on sait
déjà \deleted{\unsure{que|par}}(5.8.9) que $u_{!}(TC)\subset TC'$,

